\begin{afterword}
\summaryhead

The post extends the previous one's setup to two photons. One comes from B and one from C, and both reach the half-silvered mirror D at the same time. Using the rule that deflection multiplies the amplitude by $i$ and going straight multiplies it by 1, the author follows four routes. Two of them, both photons deflected and both going straight, end in the same configuration, one photon at each detector, with amplitudes $1$ and $-1$. These cancel. So the detectors never fire one each; both photons go to one detector or both to the other, equally often.

This works only because a configuration records ``a photon here, a photon there'' and not which photon is which. If configurations tracked identities, one-each outcomes would occur half the time. Since squared moduli do not add as probabilities do, configurations cannot be merged or split at will, as possible worlds can. The author concludes that configurations are not states of knowledge, and that physics cannot be divided into particles with individual identities.

\responsehead

The physics in this post is right. The worked amplitudes are correct, the result does not depend on the post's phase convention, and the prediction is the Hong--Ou--Mandel effect, reported in 1987, to which the post links. The central lesson, that identical particles carry no labels and that the routes into a shared configuration must be added before squaring, is standard, and the post shows it with one short calculation.

A reader learning quantum mechanics from the post should know three things it leaves out. First, the rule is stated for electrons as well as photons, but for electrons the two routes add with opposite signs, so two electrons in this experiment always leave by different sides; this was observed in 2013. Second, kind and place are not enough to make two configurations the same: the photons must also match in polarization, frequency and timing, which the next post's wording, ``a different state,'' allows for. Third, the bookkeeping drops a factor of $\sqrt{2}$ on each two-photon amplitude. That does no harm here, because the factor is the same for both detectors, but the same bookkeeping gives wrong ratios for a mirror that is not 50:50.

The post's one step beyond the physics comes near the end: ``If configurations were just states of knowledge, you could reorganize them however you liked.'' The experiment rules out photons with labels and ordinary probabilities; whether amplitudes are physical is a question of interpretation, discussed in the notes on ``Configurations and Amplitude.''

The claim that a linear rule in place of the squared modulus would make physics classical is offered without argument, but it agrees with a known result: for states with no negative entries, the evolutions that preserve the sum of the entries are those of ordinary probability.

In short: a correct account of two-photon interference, whose rule for electrons omits a sign, and whose conclusion that configurations are not knowledge goes beyond what its experiment decides.
\end{afterword}
